% -----------------------------------------------------------------------------
% Aggregated confusion-matrix figure (Full-S1 vs K1).
%
% Asset: figures/fig_confusion_aggregated_fulls1_vs_k1.pdf
%   natural size 7.16in x 8.71in, vector, generated by
%   scripts/generate_confusion_figure.py in the publication repository.
%   Provenance and verification:
%   figures/confusion_matrices/confusion_matrix_aggregation_manifest.md
%
% A plain `figure' float is used: the manuscript is single-column, so figure*
% gains nothing here and a separate float class would not be guaranteed to keep
% its numbering position relative to the other figures. Switch to figure* if the
% manuscript is migrated to a two-column journal class.
% keepaspectratio with a height cap leaves room for the caption and lets the
% float shrink for classes with a shorter text block; it does not bind here.
% -----------------------------------------------------------------------------
\begin{figure}[p]
  \centering
  \includegraphics[width=\textwidth,height=0.78\textheight,keepaspectratio]%
    {figures/fig_confusion_aggregated_fulls1_vs_k1.pdf}
  \caption{Aggregated class-level confusion matrices for Full-S1 (left column)
  and K1 (right column) on the four evaluation datasets: CWRU, JNU, HIT and
  MaFaulDa. For each dataset and model, the raw confusion counts of the nine
  matched fold--seed evaluation cells were summed element-wise, and the pooled
  matrix was then normalized by true class; each row therefore shows the
  distribution of predictions conditional on the true class, and all panels
  share the same $0$--$100\%$ colour scale; exactly-zero cells are left blank.
  The pooled matrices are descriptive summaries of the nine matched model
  evaluations rather than a single independent test split; all inferential
  comparisons rest on the paired fold--seed analysis reported in
  Section~\ref{subsec:ni_results}.}
  \label{fig:confusion_aggregated}
\end{figure}
